% ROZDZIAŁ 21 — RIESZ–MARKOW–KAKUTANI
\chapter{Twierdzenie Riesza--Markowa--Kakutaniego}
\label{ch:riesz-markov-21}

\section{Sformułowanie i rola twierdzenia}
Niech $X$ będzie lokalnie zwartą przestrzenią Hausdorffa,
a $\mathbb K\in\{\mathbb R,\mathbb C\}$.
Przestrzeń $C_c(X;\mathbb K)$ i pojęcie zwartego nośnika
wprowadzono w rozdziale~\ref{ch:radon-measures-20}.
Przez $C_0(X;\mathbb K)$ oznaczamy przestrzeń funkcji ciągłych
\emph{znikających w nieskończoności}, to znaczy takich, że
dla każdego $\varepsilon>0$ zbiór
$\{x:|f(x)|\ge\varepsilon\}$ jest zwarty.
Przestrzeń $C_0(X;\mathbb K)$ z normą supremum jest przestrzenią
Banacha. W rozważanej dalej klasie lokalnie zwartych
ośrodkowych przestrzeni metrycznych udowodnimy gęstość
$C_c(X;\mathbb K)$ w $C_0(X;\mathbb K)$.

\begin{twierdzenie}[Reprezentacja dodatnich funkcjonałów]
\label{thm:riesz-positive-21}
Niech $X$ będzie \emph{lokalnie zwartą, ośrodkową przestrzenią
metryczną}, a $L:C_c(X;\mathbb R)\to\mathbb R$ funkcjonałem
liniowym takim, że $L(f)\ge0$ dla $f\ge0$.
Istnieje dokładnie jedna miara Radona $\mu$ na $X$ spełniająca
\begin{equation}\label{eq:riesz-representation-21}
 L(f)=\int_X f(x)\,\mathrm d\mu(x),\qquad f\in C_c(X;\mathbb R).
\end{equation}
Przy czym $L$ nie musi być ograniczony w normie supremum na
całej przestrzeni $C_c(X)$.
\end{twierdzenie}

\begin{uwaga}[Zakres wersji twierdzenia]
Klasyczne twierdzenie zachodzi na ogólnych lokalnie zwartych
przestrzeniach Hausdorffa, przy właściwie określonej regularności
miary. Aby wszystkie kroki dowodu, w tym konstrukcję i regularność,
wyprowadzić tutaj z wcześniejszych rozdziałów bez ukrytych
argumentów topologicznych, udowodnimy je dla lokalnie zwartych
ośrodkowych przestrzeni metrycznych. Obejmuje to w szczególności
$\mathbb R^d$, rozmaitości drugiej przeliczalności i przestrzenie
polskie lokalnie zwarte. Wersja dla dowolnych lokalnie zwartych
przestrzeni Hausdorffa wymaga dodatkowej pracy z regularnymi
miarami zewnętrznymi; nie utożsamiamy tych dwóch zakresów.
\end{uwaga}

\section{Lemat o metrycznych miarach zewnętrznych}
\begin{definicja}[Miara zewnętrzna metryczna]
\label{def:metric-outer-21}
Niech $(Z,d)$ będzie przestrzenią metryczną.
Miara zewnętrzna $m^*:\mathcal P(Z)\to[0,\infty]$
jest \emph{metryczna}, gdy dla dowolnych $A,B\subseteq Z$
o dodatniej odległości $d(A,B)>0$ zachodzi
$m^*(A\cup B)=m^*(A)+m^*(B)$.
Dla zbioru pustego lub dowolnego niepustego przyjmujemy
$d(\varnothing,B)=+\infty$.
\end{definicja}

\begin{lemat}[Borelowska mierzalność]
\label{lem:metric-outer-borel-21}
Jeżeli $m^*$ jest metryczną miarą zewnętrzną na $(Z,d)$,
to każdy zbiór borelowski w $Z$ jest mierzalny w sensie
Carathéodory'ego.
\end{lemat}
\begin{proof}
Wystarczy wykazać mierzalność zbiorów domkniętych, gdyż
rodzina zbiorów mierzalnych Carathéodory'ego jest
$\sigma$-algebrą (rozdział~\ref{ch:caratheodory}).
Niech $F\subseteq Z$ będzie domknięty i $E\subseteq Z$.
Musimy dowieść
$m^*(E)\ge m^*(E\cap F)+m^*(E\setminus F)$;
nierówność odwrotna wynika z podaddytywności.
Jeśli $m^*(E)=+\infty$, teza jest oczywista. Załóżmy
$m^*(E)<\infty$ i oznaczmy $r(x)=d(x,F)$.
Dla każdego $n\ge1$ zbiory $E\cap F$ oraz
$E\cap\{r>1/n\}$ mają dodatnią odległość, gdyż
$r$ jest $1$-lipschitzowska. Zatem
\begin{equation}\label{eq:metric-outer-sep-21}
 m^*(E)\ge m^*(E\cap F)+m^*(E\cap\{r>1/n\}).
\end{equation}
Wykażemy, że drugi składnik po prawej zbiega do
$m^*(E\setminus F)$; nie wolno tutaj korzystać z ciągłości
z dołu, ponieważ $m^*$ jest dopiero miarą zewnętrzną.
Niech $A_0=E\cap\{r>1\}$ oraz
$A_j=E\cap\{1/(j+1)<r\le1/j\}$, $j\ge1$.
Dla ustalonej parzystości indeksów dowolne dwa
\emph{różne} zbiory $A_j$ mają dodatnią odległość.
Dla każdej skończonej rodziny takich zbiorów metryczność
i indukcja dają
$\sum m^*(A_j)=m^*(\bigcup A_j)\le m^*(E)$.
W konsekwencji obie serie, po indeksach parzystych
i nieparzystych, są zbieżne; zatem
$\sum_{j\ge1}m^*(A_j)<\infty$. Ponieważ
\[
 E\cap\{0<r\le1/n\}\subseteq\bigcup_{j\ge n}A_j,
\]
miara zewnętrzna części po lewej dąży do zera.
Z podaddytywności
\[
 m^*(E\setminus F)\le
 m^*(E\cap\{r>1/n\})+
 m^*(E\cap\{0<r\le1/n\}),
\]
a z monotoniczności wynika przeciwna nierówność
$m^*(E\cap\{r>1/n\})\le m^*(E\setminus F)$.
Przechodząc do granicy w \eqref{eq:metric-outer-sep-21},
otrzymujemy żądane kryterium Carathéodory'ego.
\end{proof}

\section{Przypadek zwarty: konstrukcja od podstaw}
\begin{twierdzenie}[Reprezentacja na zwartej przestrzeni metrycznej]
\label{thm:compact-riesz-21}
Niech $K$ będzie niepustą zwartą przestrzenią metryczną,
a $L:C(K;\mathbb R)\to\mathbb R$ dodatnim funkcjonałem
liniowym. Istnieje dokładnie jedna skończona regularna
miara borelowska $\mu$ na $K$ taka, że
$L(f)=\int_Kf\,\mathrm d\mu$ dla wszystkich $f\in C(K)$.
Ponadto $\mu(K)=L(\chi_K)$ i
$\|L\|=L(\chi_K)$.
\end{twierdzenie}
\begin{proof}
\emph{Krok 1: lokalne ograniczenie funkcjonału.}
Z $-\|f\|_\infty\chi_K\le f\le\|f\|_\infty\chi_K$
wynika $|L(f)|\le L(\chi_K)\|f\|_\infty$.
Dla $f=\chi_K$ mamy równość w normie.

\emph{Krok 2: przypisanie liczb zbiorom otwartym.}
Dla zbioru otwartego $U\subseteq K$ definiujemy
\begin{equation}\label{eq:open-content-21}
 a(U)=\sup\{L(f): f\in C(K),\ 0\le f\le1,\
          \operatorname{supp}f\subseteq U\}.
\end{equation}
W szczególności $a(\varnothing)=0$,
$0\le a(U)\le L(\chi_K)$, $a(K)=L(\chi_K)$,
a $a$ jest monotoniczna. Jeżeli $U,V$ są rozłączne
i otwarte, to $a(U\cup V)\ge a(U)+a(V)$:
wybieramy funkcje $f,g$ dopuszczalne odpowiednio dla
$U,V$; ponieważ ich nośniki są rozłączne,
$0\le f+g\le1$ i $L(f+g)=L(f)+L(g)$.

\emph{Krok 3: przeliczalna podaddytywność na otwartych.}
Jeżeli $U\subseteq\bigcup_{j\ge1}U_j$ i $f$ jest
dopuszczalna w \eqref{eq:open-content-21}, to
$\operatorname{supp}f$ ma skończone pokrycie
$U_{j_1},\ldots,U_{j_N}$. Istnieją ciągłe funkcje
$\theta_i\ge0$ o nośnikach zawartych w $U_{j_i}$,
takie, że $\sum_i\theta_i=1$ na
$\operatorname{supp}f$. Aby je zbudować, używamy
funkcji odcinających równych $1$ na skończonym
pokryciu mniejszymi zbiorami i dzielimy ich sumę
w otoczeniu nośnika; poza tym otoczeniem ustawiamy
funkcje na zero. Każde $f\theta_i$ jest nieujemne,
nie przekracza $1$ i ma nośnik w $U_{j_i}$.
Wobec tego
$L(f)=\sum_i L(f\theta_i)\le\sum_{j\ge1}a(U_j)$.
Po wzięciu supremum otrzymujemy
$a(U)\le\sum_j a(U_j)$.

\emph{Krok 4: miara zewnętrzna.}
Dla dowolnego $A\subseteq K$ definiujemy
\begin{equation}\label{eq:riesz-outer-21}
 m^*(A)=\inf\{a(U):A\subseteq U,\ U\text{ otwarty}\}.
\end{equation}
Jeżeli $A\subseteq\bigcup_j A_j$, dobieramy otwarte
$U_j\supseteq A_j$ z
$a(U_j)\le m^*(A_j)+\varepsilon2^{-j}$ (przyjmujemy
osobno trywialny przypadek nieskończonej sumy).
Krok 3 daje
$m^*(A)\le a(\bigcup_jU_j)\le\sum_jm^*(A_j)+\varepsilon$.
Zatem $m^*$ jest miarą zewnętrzną.
Dla zbioru otwartego $U$ mamy $m^*(U)=a(U)$, ponieważ
$U\subseteq V$ implikuje $a(U)\le a(V)$.
Jeśli $d(A,B)=\delta>0$, to dla każdego otwartego
$W\supseteq A\cup B$ zbiory
$W_A=W\cap\{x:d(x,A)<\delta/3\}$ i
$W_B=W\cap\{x:d(x,B)<\delta/3\}$ są otwarte,
rozłączne i zawierają odpowiednio $A,B$. Stąd
$a(W)\ge a(W_A)+a(W_B)\ge m^*(A)+m^*(B)$.
Biorąc infimum po $W$ otrzymujemy metryczność $m^*$.
Lemat~\ref{lem:metric-outer-borel-21} zapewnia, że
wszystkie zbiory borelowskie są mierzalne. Ograniczenie
$\mu=m^*|_{\mathcal B(K)}$ jest więc miarą, a
$\mu(K)=L(\chi_K)<\infty$.

\emph{Krok 5: regularność.}
Regularność zewnętrzna wynika z \eqref{eq:riesz-outer-21}.
Dla otwartego $U$ i $\varepsilon>0$ dobieramy
$0\le f\le1$ z $\operatorname{supp}f\subseteq U$ i
$L(f)>\mu(U)-\varepsilon$. Przyjmijmy
$H=\operatorname{supp}f$, który jest zwarty.
Każdy otwarty $V\supseteq H$ spełnia
$L(f)\le a(V)$, zatem
$\mu(H)=\inf_{V\supseteq H}a(V)\ge L(f)$.
To dowodzi regularności wewnętrznej dla otwartych.
Dla dowolnego borelowskiego $E$ i $\varepsilon>0$
wybieramy otwarty $U\supseteq E$ z
$\mu(U\setminus E)<\varepsilon$.
Następnie wybieramy zwarty $H\subseteq U$ z
$\mu(U\setminus H)<\varepsilon$ oraz otwarty
$W\supseteq U\setminus E$ z
$\mu(W)<\mu(U\setminus E)+\varepsilon$.
Zbiór $H\setminus W$ jest zwarty, zawiera się w $E$,
a $\mu(E\setminus(H\setminus W))<3\varepsilon$.
Otrzymujemy regularność wewnętrzną względem zwartych.

\emph{Krok 6: identyfikacja całki z $L$.}
Niech $f\in C(K)$, $f\ge0$, $M=\|f\|_\infty>0$.
Dla $N\ge1$ połóżmy $\delta=M/N$ i
$U_i=\{f>i\delta\}$, $i=0,\ldots,N$.
Dla $i=0,\ldots,N-1$ zbiór
$H_i=\{f\ge(i+1)\delta\}$ jest zwarty i zawiera się
w $U_i$. Istnieje $h_i\in C(K)$, $0\le h_i\le1$,
$h_i=1$ na $H_i$, $\operatorname{supp}h_i\subseteq U_i$;
wynika to z lematu o funkcji odcinającej
\ref{lem:cutoff-radon-20}. Punktowo
$f\le\delta\chi_K+\delta\sum_{i=0}^{N-1}h_i$.
Z dodatniości $L$ dostajemy
\begin{equation}\label{eq:riesz-upper-sum-21}
 L(f)\le\delta L(\chi_K)+
           \delta\sum_{i=0}^{N-1}\mu(U_i).
\end{equation}
Z drugiej strony dla dowolnych funkcji
$g_i\in C(K)$, $0\le g_i\le1$,
$\operatorname{supp}g_i\subseteq U_i$ ($1\le i\le N$),
zachodzi $\delta\sum_{i=1}^{N}g_i\le f$ punktowo.
Zatem, biorąc niezależnie supremum po $g_i$,
\begin{equation}\label{eq:riesz-lower-sum-21}
 \delta\sum_{i=1}^{N}\mu(U_i)\le L(f).
\end{equation}
Natomiast rozcięcie poziomów funkcji daje
\[
 \delta\sum_{i=1}^{N}\mu(U_i)
 \le\int_Kf(x)\,\mathrm d\mu(x)
 \le\delta\sum_{i=0}^{N-1}\mu(U_i).
\]
Różnica między lewą i prawą sumą wynosi
$\delta\mu(U_0)\le\delta L(\chi_K)$.
Z \eqref{eq:riesz-upper-sum-21} i
\eqref{eq:riesz-lower-sum-21} wynika
$|L(f)-\int_K f\,\mathrm d\mu|\le2\delta L(\chi_K)$.
Wysyłając $N\to\infty$, otrzymujemy reprezentację
$f\ge0$. Ogólny $f$ rozkładamy na części dodatnią
i ujemną.

\emph{Krok 7: jednoznaczność.}
Dla dowolnej miary regularnej $\rho$ reprezentującej $L$
i otwartego $U$ regularność oraz funkcje odcinające
dają $\rho(U)=a(U)$ (por.
\eqref{eq:cc-open-sup-20}). Stąd z regularności
zewnętrznej $\rho=\mu$ na wszystkich zbiorach
borelowskich.
\end{proof}

\section{Przestrzenie lokalnie zwarte: lokalizacja i sklejanie}
\begin{lemat}[Wyczerpanie przez zbiory relatywnie zwarte]
\label{lem:exhaustion-21}
Niech $X$ będzie lokalnie zwartą ośrodkową przestrzenią
metryczną. Istnieją otwarte zbiory $U_n$ takie, że
$U_n\subseteq\overline U_n\subseteq U_{n+1}$,
$\overline U_n$ jest zwarty i $\bigcup_nU_n=X$.
Ponadto istnieją $\chi_n\in C_c(X)$ spełniające
$0\le\chi_n\le1$, $\chi_n=1$ na $\overline U_n$
oraz $\operatorname{supp}\chi_n\subseteq U_{n+1}$.
\end{lemat}
\begin{proof}
Lokalna zwartość daje każdemu punktowi relatywnie zwarte
otwarte otoczenie. Ośrodkowa przestrzeń metryczna jest
drugiej przeliczalności, więc z tego pokrycia wybieramy
przeliczalne podpokrycie $(V_j)_{j\ge1}$.
Indukcyjnie dobieramy relatywnie zwarte otwarte $U_n$,
obejmujące $\overline U_{n-1}\cup\overline V_1\cup\cdots
\cup\overline V_n$: skończone pokrycie tego zwartego
zbioru relatywnie zwartymi otoczeniami ma relatywnie
zwartą sumę. Dobierając otwarte otoczenie o zwartej
domkniętej powłoce, uzyskujemy żądane zawierania.
Funkcje $\chi_n$ istnieją na mocy lematu
\ref{lem:cutoff-radon-20}.
\end{proof}

\begin{proof}[Dowód twierdzenia~\ref{thm:riesz-positive-21}]
\emph{Krok 1: dodatniość lokalna i zwartość.}
Dla dowolnego zwartego $H\subseteq X$ wybieramy
$\eta\in C_c(X)$, $\eta\ge0$, $\eta=1$ na $H$.
Jeśli $\operatorname{supp}f\subseteq H$, to
$-\|f\|_\infty\eta\le f\le\|f\|_\infty\eta$,
więc $|L(f)|\le\|f\|_\infty L(\eta)$.
Wybierzmy wyczerpanie i funkcje $\chi_n$ z lematu
\ref{lem:exhaustion-21}; oznaczmy $K_n=\overline U_n$
i $J_n=\overline U_{n+1}$.
Dla $h\in C(J_n)$ funkcję $\chi_nh$ na $J_n$
przedłużamy przez zero na $X\setminus J_n$.
Jest ona ciągła, ponieważ
$\operatorname{supp}\chi_n\subseteq U_{n+1}$.
Definiujemy $L_n(h)=L(\chi_nh)$; jest to dodatni
funkcjonał na $C(J_n)$.
Z twierdzenia~\ref{thm:compact-riesz-21}
istnieje skończona regularna miara $\rho_n$
na $J_n$ reprezentująca $L_n$. Przedłużamy ją
przez zero na dopełnienie $J_n$, otrzymując miarę
borelowską na $X$ (także oznaczaną $\rho_n$).

\emph{Krok 2: zgodność miar na wyczerpaniu.}
Jeżeli $m\ge n$, to $\chi_m=\chi_n=1$ na $U_n$.
Dla dowolnego $f\in C_c(U_n)$ zachodzi zatem
\[
 \int_Xf\,\mathrm d\rho_m=L(f)
  =\int_Xf\,\mathrm d\rho_n.
\]
Ograniczenia $\rho_m|_{U_n}$ i $\rho_n|_{U_n}$
są miarami Radona (na otwartej przestrzeni $U_n$).
Jednoznaczność z twierdzenia~\ref{thm:uniqueness-cc-20}
daje $\rho_m|_{U_n}=\rho_n|_{U_n}$.
Uwaga: ograniczenie miary Radona do otwartego podzbioru
jest Radona; regularność wewnętrzna oraz zewnętrzna
wynika z regularności wyjściowej i przecięcia z podprzestrzenią.

\emph{Krok 3: sklejenie.}
Stawiamy $U_0=\varnothing$, $B_n=U_n\setminus U_{n-1}$.
Dla borelowskiego $E\subseteq X$ określamy
\begin{equation}\label{eq:gluing-measures-21}
 \mu(E)=\sum_{n\ge1}\rho_n(E\cap B_n).
\end{equation}
Ponieważ $B_n$ są parami rozłączne, a suma miar jest
miarą, $\mu$ jest miarą borelowską. Dla $E\subseteq U_m$
przeliczalna addytywność i krok 2 dają
$\mu(E)=\rho_m(E)$; dotyczy to również wszystkich
$E\in\mathcal B(U_m)$. W szczególności
$\mu(U_m)=\rho_m(U_m)<\infty$.
Zatem $\mu$ jest lokalnie skończona. Rozdział
\ref{ch:radon-measures-20} pokazuje, że na lokalnie
zwartej ośrodkowej przestrzeni metrycznej jest ona
miarą Radona.

\emph{Krok 4: reprezentacja.}
Jeśli $f\in C_c(X)$, jej nośnik mieści się w pewnym
$U_m$ (ponieważ $(U_m)$ jest rosnącym pokryciem
zwartego nośnika). Wtedy
\[
 \int_X f(x)\,\mathrm d\mu(x)
 =\int_X f(x)\,\mathrm d\rho_m(x)
 =L(\chi_m f)=L(f).
\]
Jednoznaczność wynika z twierdzenia
\ref{thm:uniqueness-cc-20}.
\end{proof}

\section{Przestrzeń $C_0(X)$ oraz miary o skończonej wariacji}
\begin{lemat}[Gęstość $C_c(X)$ w $C_0(X)$]
\label{lem:density-cc-c0-21}
Dla lokalnie zwartej ośrodkowej przestrzeni metrycznej
$X$ przestrzeń $C_c(X;\mathbb K)$ jest gęsta w
$C_0(X;\mathbb K)$ w normie supremum.
\end{lemat}
\begin{proof}
Niech $f\in C_0(X)$ i $\varepsilon>0$.
Zbiór $H=\{|f|\ge\varepsilon\}$ jest zwarty.
Wybieramy $\eta\in C_c(X)$ równą $1$ na $H$,
$0\le\eta\le1$. Wówczas $\eta f\in C_c(X)$
i $\|f-\eta f\|_\infty\le\varepsilon$.
\end{proof}

\begin{twierdzenie}[Ograniczone dodatnie funkcjonały na $C_0$]
\label{thm:riesz-c0-positive-21}
W założeniach twierdzenia~\ref{thm:riesz-positive-21}
niech $L$ będzie ograniczonym dodatnim funkcjonałem
na $C_0(X;\mathbb R)$. Wtedy jego miara reprezentująca
jest skończona i $\|L\|=\mu(X)$. Odwrotnie każda
skończona miara Radona definiuje taki funkcjonał.
\end{twierdzenie}
\begin{proof}
Dla kolejnych zbiorów $U_n$ z lematu
\ref{lem:exhaustion-21} wybieramy $\chi_n=1$
na $\overline U_n$, $0\le\chi_n\le1$.
Wówczas $\mu(U_n)\le\int_X\chi_n\,\mathrm d\mu
=L(\chi_n)\le\|L\|$. Przechodząc do granicy,
$\mu(X)\le\|L\|$. W przeciwną stronę,
$|L(f)|\le\mu(X)\|f\|_\infty$ dla $f\in C_c$;
przez gęstość nierówność zachodzi na $C_0$.
Każda skończona miara Radona definiuje przez całkę
funkcjonał liniowy ciągły i dodatni, co kończy dowód.
\end{proof}

\begin{definicja}[Skończona zespolona miara Radona]\label{idx:21:definicja:08}

Niech $X$ będzie lokalnie zwartą ośrodkową przestrzenią
metryczną. Skończoną zespoloną miarę borelowską $\nu$
nazywamy miarą Radona, jeśli jej wariacja całkowita
$|\nu|$ jest skończoną miarą Radona. Wówczas części
rzeczywista i urojona są miarami ze znakiem o skończonej
wariacji, a wszystkie całki z funkcji $C_0(X)$ istnieją.
\end{definicja}

\begin{twierdzenie}[Reprezentacja zespolona i norma wariacji]
\label{thm:riesz-complex-21}
Niech $X$ będzie lokalnie zwartą ośrodkową przestrzenią
metryczną, a $T:C_0(X;\mathbb C)\to\mathbb C$
ograniczonym funkcjonałem liniowym. Istnieje dokładnie
jedna skończona zespolona miara Radona $\nu$ taka, że
\[
 T(f)=\int_X f(x)\,\mathrm d\nu(x),\qquad f\in C_0(X;\mathbb C).
\]
Ponadto $\|T\|=|\nu|(X)$, gdzie $|\nu|$ oznacza
wariację całkowitą miary zespolonej.
\end{twierdzenie}
\begin{proof}
Dla nieujemnego $g\in C_c(X;\mathbb R)$ zdefiniujmy
\[
 M(g)=\sup\{|T(h)|:h\in C_c(X;\mathbb C),\ |h|\le g\}.
\]
Jest to skończone, bo $M(g)\le\|T\|\|g\|_\infty$.
Ma własności $M(ag)=aM(g)$ dla $a\ge0$ oraz
$M(g_1+g_2)=M(g_1)+M(g_2)$.
Dla nierówności „$\le$” rozkładamy dowolną $h$,
$|h|\le g_1+g_2$, na $h_i=h g_i/(g_1+g_2)$ na
$\{g_1+g_2>0\}$, a w pozostałych punktach przyjmujemy
$h_i=0$. Ciągłość przy zerach mianownika wynika
z $|h_i|\le g_i$; mamy $|h_i|\le g_i$ i
$h_1+h_2=h$. Dla nierówności „$\ge$” bierzemy
$h_i$ niemal realizujące supremum i mnożymy je przez
stałe fazy tak, aby wartości $T(h_i)$ były nieujemne
rzeczywiste. Wtedy $|h_1+h_2|\le g_1+g_2$ i
$T(h_1+h_2)\approx M(g_1)+M(g_2)$.
Z addytywności i dodatniej jednorodności $M$ przedłuża
się jednoznacznie do dodatniego rzeczywistego funkcjonału
na $C_c(X;\mathbb R)$: $M(f)=M(f^+)-M(f^-)$.
Ponieważ $M(g)\le\|T\|\|g\|_\infty$, przedłużony
funkcjonał $M$ jest ograniczony w normie supremum
i rozciąga się przez gęstość na $C_0(X;\mathbb R)$.
Twierdzenie~\ref{thm:riesz-c0-positive-21} daje
skończoną miarę Radona $\sigma$ taką, że
$M(g)=\int_Xg\,\mathrm d\sigma$ dla $g\ge0$.
W szczególności dla $h\in C_c(X;\mathbb C)$ zachodzi
\begin{equation}\label{eq:riesz-domination-21}
 |T(h)|\le \int_X|h(x)|\,\mathrm d\sigma(x).
\end{equation}
Przestrzeń $C_c(X;\mathbb C)$ jest gęsta w
$L^1(\sigma;\mathbb C)$: przybliżamy funkcje charakterystyczne
zbiorów borelowskich przez funkcje odcinające między
zwartym podzbiorem i otwartym nadzbiorem o małej
różnicy miar, a następnie przybliżamy funkcje proste.
Zatem $T$ jednoznacznie przedłuża się do funkcjonału
na $L^1(\sigma)$ o normie co najwyżej $1$.
Dualność $L^1$--$L^\infty$ dla skończonej miary
(\ref{ch:dualnosc-lp}) daje $w\in L^\infty(\sigma)$,
$|w|\le1$ $\sigma$-prawie wszędzie, takie że
$T(h)=\int_X h w\,\mathrm d\sigma$.
Kładziemy $\nu=w\sigma$; jest to skończona zespolona
miara Radona, ponieważ $|\nu|\le\sigma$ i
regularność przechodzi na miary zdominowane.
Reprezentacja obowiązuje dla $C_c$, a przez gęstość
także dla $C_0$.

Pozostaje identyczność norm. Nierówność
$\|T\|\le|\nu|(X)$ jest bezpośrednia. Dla odwrotnej
wybieramy mierzalną funkcję fazową $u$ o $|u|\le1$
taką, że $u\,\mathrm d\nu=\mathrm d|\nu|$.
Z regularności $|\nu|$ i aproksymacji funkcji
mierzalnych przez $C_c$ w $L^1(|\nu|)$ istnieją
$f_n\in C_c(X;\mathbb C)$, $\|f_n\|_\infty\le1$,
$f_n\to u$ w $L^1(|\nu|)$. Ograniczenie normy
zapewnia obcięcie aproksymantów do domkniętego
dysku jednostkowego, które nie zwiększa błędu
punktowego względem $u$. Stąd
$T(f_n)\to\int_Xu\,\mathrm d\nu=|\nu|(X)$,
więc $\|T\|\ge|\nu|(X)$.
Jednoznaczność $\nu$ wynika z jednoznaczności dla
części rzeczywistej i urojonej po rozdzieleniu miar
oraz z regularności; szczegółowo równość całek na
$C_c$ implikuje równość każdej części miar ze znakiem,
a to wynika z twierdzenia~\ref{thm:uniqueness-cc-20}
po zastosowaniu rozkładu Jordana.
\end{proof}

\begin{przyklad}[Funkcjonał ewaluacji]
Dla ustalonego $x_0\in X$ funkcjonał $T(f)=f(x_0)$
na $C_0(X)$ jest ciągły i ma normę $1$.
Jego miarą reprezentującą jest $\delta_{x_0}$.
\end{przyklad}

\begin{przyklad}[Miara nieskończona i funkcjonał lokalny]
Dla $X=\mathbb R$ i $L(f)=\int_{\mathbb R} f(x)
\,\mathrm d\lambda_1(x)$, $f\in C_c(\mathbb R)$,
otrzymujemy dodatni funkcjonał nieograniczony w normie
supremum na całym $C_c$: funkcje odcinające równe $1$
na $[-n,n]$ mają normę supremum $1$, ale całki
co najmniej $2n$. Dlatego $L$ nie przedłuża się
ciągle do $C_0(\mathbb R)$.
\end{przyklad}

\begin{wniosek}[Porządek dodatnich funkcjonałów i miar]
\label{cor:riesz-order-21}
Niech $X$ będzie lokalnie zwartą ośrodkową przestrzenią
metryczną, a $L_1,L_2:C_c(X;\mathbb R)\to\mathbb R$
dodatnimi funkcjonałami liniowymi, reprezentowanymi
odpowiednio przez miary Radona $\mu_1,\mu_2$.
Wówczas następujące warunki są równoważne:
\begin{enumerate}[label=(\roman*)]
\item $L_1(f)\le L_2(f)$ dla każdego nieujemnego $f\in C_c(X)$;
\item $\mu_1(E)\le\mu_2(E)$ dla każdego zbioru borelowskiego
$E\subseteq X$.
\end{enumerate}
\end{wniosek}
\begin{proof}
Implikacja (ii)$\Rightarrow$(i) wynika z monotoniczności całki.
Załóżmy (i). Dla otwartego $U\subseteq X$ zachodzi
\[
 \mu_1(U)=\sup\{L_1(f):f\in C_c(X),\ 0\le f\le1,
 \ \operatorname{supp}f\subseteq U\}.
\]
Wzór ten udowodniono w rozdziale~\ref{ch:radon-measures-20};
obejmuje także przypadek $\mu_1(U)=+\infty$.
Nierówność dla funkcjonałów daje
$\mu_1(U)\le\mu_2(U)$. Dla dowolnego borelowskiego
$E$ z regularności zewnętrznej otrzymujemy
$\mu_1(E)=\inf_{U\supseteq E}\mu_1(U)
\le\inf_{U\supseteq E}\mu_2(U)=\mu_2(E)$.
\end{proof}

\begin{uwaga}[Gdzie użyliśmy twierdzenia Tichonowa?]
W konstrukcji na zwartej przestrzeni metrycznej nie
potrzebowaliśmy twierdzenia Tichonowa: miarę zewnętrzną
zbudowaliśmy bezpośrednio z dodatniego funkcjonału.
W rozdziale~\ref{ch:haar-23}, poświęconym mierze Haara,
pojawi się natomiast zwartość produktu nieprzeliczalnie
wielu przedziałów.
Nie należy jej mylić ze zwartą metryczną przestrzenią,
dla której wystarczy argument diagonalny.
\end{uwaga}

\section{Konsekwencje: oddzielenie miar i problem momentów}
\begin{wniosek}[Miary odróżniane przez funkcje ciągłe]
\label{cor:measures-separate-21}
Niech $\mu,\nu$ będą miarami Radona na lokalnie zwartej
ośrodkowej przestrzeni metrycznej $X$. Wtedy
$\mu=\nu$ wtedy i tylko wtedy, gdy
$\int_Xf\,\mathrm d\mu=\int_Xf\,\mathrm d\nu$
dla wszystkich $f\in C_c(X;\mathbb R)$.
\end{wniosek}
\begin{proof}
Implikacja w prawo jest oczywista, a odwrotna to
\ref{thm:uniqueness-cc-20}.
\end{proof}

\begin{uwaga}[Problem momentów jako ostrzeżenie]
Jeżeli na $\mathbb R$ znamy całki
$\int_{\mathbb R}x^k\,\mathrm d\mu(x)$ dla wszystkich
$k\ge0$, nie zawsze wyznacza to jednoznacznie miarę.
Funkcje wielomianowe nie należą do $C_0(\mathbb R)$
i nie ma automatycznej gęstości wielomianów w normie
supremum na całej prostej. Na zwartym przedziale
sytuacja jest inna: twierdzenie Stone'a--Weierstrassa
przybliża dowolną funkcję ciągłą wielomianami,
więc momenty jednoznacznie wyznaczają skończoną miarę.
\end{uwaga}

\section{Podsumowanie i przygotowanie do zbieżności miar}
Reprezentacja dodatnich funkcjonałów pozwala przejść
od analizy na $C_c(X)$ do teorii miar Radona. Dla
funkcjonałów na $C_0(X)$ skończoność normy odpowiada
skończoności masy, zaś dla zespolonych funkcjonałów
norma jest wariacją całkowitą. W rozdziale~\ref{ch:weak-measures-22}
zastosujemy zwarty przypadek twierdzenia do wykazania
ciągowej zwartości przestrzeni prawdopodobieństw na
zwartej przestrzeni metrycznej.
